% !TEX root = ../main.tex
% CHU SO HUU: Mỹ Uyên và Mai Quyên | KHONG THEO PREAMBLE, KHONG BEGIN DOCUMENT TRONG THAN PHAN NAY.
\documentclass[../main.tex]{subfiles}
\begin{document}
\section{Kết quả phân tích tương quan Pearson}
\label{sec:t31:ket-qua}

\subsection{Dữ liệu, ma trận tương quan và heatmap}
Tập huấn luyện gồm $125\,973$ quan sát. Tệp \nolinkurl{train_clean.csv}
có 45 cột: 41 đặc trưng gốc và bốn cột nhãn/metadata. Dữ liệu không
có giá trị thiếu hoặc bản ghi trùng. Theo quy trình ở
Mục~\ref{sec:t31:pearson-r}, còn 37 biến số không hằng (Bảng~\ref{tab:t31:data}).

\begin{table}[htbp]
\centering
\caption{Tóm tắt dữ liệu đầu vào của phân tích Pearson}
\label{tab:t31:data}
\small
\begin{tabular}{lr}
\toprule
Thông tin & Kết quả\\
\midrule
Số quan sát (tập huấn luyện) & $125\,973$\\
Số đặc trưng gốc & 41\\
Số cột nhãn/metadata trong dữ liệu sạch & 4\\
Giá trị thiếu / bản ghi trùng & 0 / 0\\
Biến hằng không đưa vào Pearson & \texttt{num\_outbound\_cmds}\\
Số biến định lượng được phân tích & 37\\
\bottomrule
\end{tabular}
\end{table}

Ma trận tương quan có kích thước $37\times37$, đối xứng, đường chéo chính
bằng 1 và không có giá trị không xác định hoặc vô cực. Số cặp khác nhau
là $37(37-1)/2=666$. Ma trận đầy đủ được lưu trong
\texttt{pearson\_matrix.csv}; danh sách các cặp vượt ngưỡng được lưu
trong \texttt{high\_correlation\_pairs.csv} (các tệp này nằm trong thư mục
kèm theo báo cáo, xem Phụ lục~\ref{app:files}). Hình~\ref{fig:t31:heatmap}
cho thấy các cụm tương quan; việc xác định cặp vượt ngưỡng được tính
trực tiếp từ hệ số chứ không suy từ màu sắc.

\begin{figure}[H]
\centering
\includegraphics[width=0.60\linewidth]{figures/task31_pearson_heatmap.png}
\caption{Heatmap Pearson của 37 đặc trưng định lượng trên tập \texttt{KDDTrain+}.}
\label{fig:t31:heatmap}
\end{figure}
\FloatBarrier

\subsection{Các cặp tương quan rất cao và khả năng dư thừa}
Lọc tam giác trên của ma trận với $|r|>0{,}95$ cho 10/666 cặp;
\emph{tất cả đều tương quan thuận} (Bảng~\ref{tab:t31:highcorr}).

Trong 10 cặp, một cặp liên quan đến biến đếm \emph{root}, sáu cặp thuộc
nhóm tỷ lệ lỗi SYN và ba cặp thuộc nhóm tỷ lệ lỗi REJ. Ở nhóm root,
\texttt{num\_compromised} đếm các điều kiện compromised, còn
\texttt{num\_root} đếm số lần truy cập root; hai đặc trưng liên hệ chặt
nhưng không đồng nghĩa về định nghĩa. Ở nhóm SYN, các biến
\nolinkurl{serror_rate}, \nolinkurl{srv_serror_rate},
\nolinkurl{dst_host_serror_rate} và \nolinkurl{dst_host_srv_serror_rate}
đo tỷ lệ lỗi tại các phạm vi tổng hợp khác nhau. Sáu cặp có
$0{,}977596\leq r\leq0{,}993289$. Ở nhóm REJ, ba cặp vượt ngưỡng
liên quan đến \nolinkurl{rerror_rate}, \nolinkurl{srv_rerror_rate} và
\nolinkurl{dst_host_srv_rerror_rate}, với
$0{,}964449\leq r\leq0{,}989008$. Những phạm vi thống kê khác nhau có
thể mang thông tin riêng dù tương quan rất cao.

\subsection{Kết quả kiểm định giả thuyết Pearson}
Bảng~\ref{tab:t31:tests} báo cáo kiểm định hai phía cho chín cặp minh hoạ,
dùng cùng $n=125\,973$ kết nối và $df=125\,971$.

Có tám trong chín cặp cho $p<0{,}05$; kết quả quyết định không đổi sau
điều chỉnh Holm cho đúng chín phép thử được trình bày. Các cặp P1--P4 có
$r>0{,}97$, phù hợp với tương quan tuyến tính thuận rất mạnh trong
Bảng~\ref{tab:t31:highcorr}. Cặp P5 (cùng tỷ lệ dịch vụ) có
$r=0{,}788978$, còn P6 (số kết nối) có $r=0{,}471079$. Cặp P7
liên hệ lỗi SYN và REJ có $r=-0{,}227536$, phản ánh tương quan
nghịch tương đối yếu trong mẫu.

Đặc biệt, P8 (\texttt{duration}, \texttt{src\_bytes}) có
$r=0{,}070737$ nhưng $p=1{,}929\cdot10^{-139}$: với cỡ mẫu rất lớn,
mối liên hệ tuyến tính rất yếu vẫn có thể đạt ý nghĩa thống kê.
Ngược lại, P9 (\texttt{src\_bytes}, \texttt{dst\_bytes}) cho
$r=0{,}000204$ và $p=0{,}9423$, nên chưa bác bỏ $H_0$;
điều đó \emph{không chứng minh} hai đặc trưng độc lập. Khoảng tin cậy
xấp xỉ 95\% theo biến đổi Fisher cho $\rho$ của P9 là
$[-0{,}005318;\,0{,}005726]$. Giá trị $p$ rất nhỏ bị làm tròn thành 0
trong tính toán số được trình bày $p<10^{-300}$, không phải xác suất bằng 0.

\subsection{Thảo luận và đề xuất đặc trưng cần xem xét}
Cặp root P1 nhạy với các quan sát có số đếm cực lớn. Để chẩn đoán,
nhóm tính lại $r$ sau khi chỉ giữ các dòng có \emph{cả hai} số đếm
không vượt 10. Có $86/125\,973$ dòng bị tạm loại khỏi phép kiểm tra độ nhạy,
và $r$ giảm từ $0{,}998833$ xuống $0{,}095623$. Đây \emph{không phải}
bước làm sạch dữ liệu: các trường hợp cực trị có thể quan trọng cho nhận diện
xâm nhập. Kết quả nhắc rằng tương quan cao chưa chắc bền vững trên mọi
phần dữ liệu.

Trên cơ sở kết quả mô tả và kiểm định, nhóm đưa năm đặc trưng vào danh sách
\emph{cần xem xét tiếp}: \nolinkurl{num_root}, \nolinkurl{srv_serror_rate},
\nolinkurl{dst_host_srv_serror_rate}, \nolinkurl{srv_rerror_rate},
\nolinkurl{dst_host_srv_rerror_rate}. Đây không phải danh sách các biến đã bị
xóa. Những tỷ lệ lỗi được biểu diễn trong khoảng $[0,1]$, không phải
phần trăm $0$--$100$. Chưa có thử nghiệm hiệu quả mô hình trước và sau khi
loại biến, nên không thể kết luận đặc trưng nào hoàn toàn thay thế được
đặc trưng khác. Quyết định chọn biến còn cần dựa vào định nghĩa mạng,
độ nhạy với ngoại lệ và đánh giá mô hình độc lập.

\end{document}
